\chapter{Классификация linking bridge внутри существующего 42-carrier}
\label{ch:v8-baryon-c0-existing-42-carrier-linking-bridge-classification-gate}

Минимальный $M_2$-мост должен соединять нейтральную исходную прямую с
нейтральным auxiliary mass carrier. Проверим, имеется ли такая стрелка в
уже построенном концевом пространстве. Его gauge-представление разлагается
как
\begin{equation}
 H_{21}=(3,2)_{1/6}\oplus(3,1)_{2/3}\oplus(3,1)_{-1/3}
 \oplus3(1,2)_{-1/2}\oplus3(1,1)_{-1}.
 \label{eq:v8-baryon-c0-endpoint-representation-ledger}
\end{equation}
Тривиального представления $(1,1)_0$ в этом списке нет, поэтому
\begin{equation}
 \operatorname{Hom}_G(\mathbb C_0,H_{21})=0.
 \label{eq:v8-baryon-c0-neutral-endpoint-hom-zero}
\end{equation}
Эквивалентно, генератор гиперзаряда не имеет нулевых диагональных весов и
удовлетворяет
\begin{equation}
 \rank Y=21,
 \qquad \ker Y=0.
 \label{eq:v8-baryon-c0-hypercharge-no-neutral-vector}
\end{equation}

\section{Опорный запрет полного кадра}

Все текущие направления являются матрицами на том же концевом носителе:
\begin{equation}
 \mathcal F_{42}=\operatorname{span}_{\mathbb R}\{F_a\}_{a=1}^{42}
 \subset\operatorname{End}(H_{21}).
 \label{eq:v8-baryon-c0-existing-frame-ambient-algebra}
\end{equation}
После явного добавления одной внешней auxiliary-линии
$H_{22}=H_{21}\oplus\mathbb C_{\rm aux}$ старый кадр вкладывается как
\begin{equation}
 \iota(F)=\begin{pmatrix}F&0\\0&0\end{pmatrix},
 \qquad
 P_{\rm aux}=\begin{pmatrix}0_{21}&0\\0&1\end{pmatrix}.
 \label{eq:v8-baryon-c0-old-frame-zero-extension}
\end{equation}
Для любого $X\in\mathcal F_{42}$ точно
\begin{equation}
 \iota(X)P_{\rm aux}=0=P_{\rm aux}\iota(X),
 \qquad
 P_{\rm aux}\iota(X)^*\iota(X)P_{\rm aux}=0.
 \label{eq:v8-baryon-c0-old-frame-annihilates-auxiliary-line}
\end{equation}
Следовательно, imprimitivity identity
\begin{equation}
 E^*E=P_{\rm aux}
 \label{eq:v8-baryon-c0-required-imprimitivity-identity}
\end{equation}
не имеет решения внутри старого 42-мерного span.

\section{Ближайшие нецелевые обходы}

Если auxiliary-линии назначить заряд $-1$, одномерное представление уже
встречается, но с кратностью три:
\begin{equation}
 \dim\operatorname{Hom}_G((1,1)_{-1},H_{21})=3.
 \label{eq:v8-baryon-c0-charged-singlet-hom-dimension}
\end{equation}
Такой мост не нейтрален и не уникален. Ранее найденный единственный
Real-поддержанный канал также не подходит: он имеет тип
\begin{equation}
 Y_R\longrightarrow Q_L\rightsquigarrow u_R,
 \qquad \rank E_{\rm old}=3,
 \label{eq:v8-baryon-c0-old-real-channel-rank}
\end{equation}
то есть является цветовым триплетным каналом, а не скалярной матричной
единицей между двумя нейтральными прямыми.

\begin{theorem}[Existing-carrier bridge no-go]
Текущий 42-мерный кадр не содержит требуемого linking bridge. Во-первых,
он действует только на $H_{21}$ и аннулирует добавленную auxiliary-линию.
Во-вторых, $H_{21}$ не содержит нейтрального gauge-синглета, который мог бы
служить исходным endpoint углом. Ближайшие заряженные или цветные каналы
имеют неверный тип либо неединственную кратность.
\end{theorem}

\begin{nogocriterion}[Оболочка не создаёт отсутствующий endpoint]
Нельзя получить $E^*E=P_{\rm aux}$ линейными комбинациями, произведениями
или замыканием старых операторов, если вся исходная алгебра действует
нулём на $\mathbb C_{\rm aux}$. Необходим новый endpoint state и новая
off-diagonal стрелка.
\end{nogocriterion}

\begin{protocol}\begin{enumerate}
\item нейтральные векторы в $H_{21}$: \textbf{$0$};
\item старый frame действует на auxiliary line: \textbf{нет};
\item решений $E^*E=P_{\rm aux}$ в старом span: \textbf{$0$};
\item заряженный singlet обход: \textbf{кратность $3$, неверный тип};
\item старый Real-канал: \textbf{ранг $3$, цветной};
\item следующий гейт: \textbf{минимальное нейтральное endpoint-расширение}.
\end{enumerate}\end{protocol}

Машинный сертификат сохранён в
\path{s2t_v8_baryon_c0_existing_42_carrier_linking_bridge_classification_gate_results.json}.